\documentclass{situationist-preprint}

\title{The Best System and the Problem of Simplicity}
\msid{MET-12-b}
\msfamily{MET-12}
\msdiscipline{Metaphysics}
\mstype{Companion manuscript}
\mskeywords{laws of nature; best system account; simplicity; strength; vagueness; natural properties; robustness}

\begin{document}
\maketitle

\begin{abstract}
On the best system account defended in MET-12-a, the laws are the axioms of the true system that best balances simplicity and strength. The account is often charged with making lawhood depend on standards that are vague, language-relative or arbitrary. We separate three versions of the charge. Language relativity is removed by restricting the vocabulary of systems to genuine properties and relations. The weighting of simplicity against strength is vague, but vagueness in the standards yields fewer laws, not arbitrary ones: a generalization is a law only if it belongs to the best systems on every admissible weighting. Finally, the dependence on the standards the sciences actually use is not a dependence on us, since it fixes only which summary of a mind-independent mosaic counts as best.
\end{abstract}

\section{Three versions of the charge}

MET-12-a (Definitions~2.2--2.4) defines a system as a set of axioms in a language whose primitive predicates express fundamental natural properties, its strength as the information it conveys about the mosaic and its simplicity as inversely measured by the number and complexity of its axioms; the fundamental laws are the axioms of the best system, and a generalization is a law if every best system entails it. Objection~3 of MET-12-a raises the problem of simplicity and defers its detailed treatment. The problem has three versions.

\emph{Language relativity.} Simplicity, measured by the length of axioms, depends on the vocabulary. With a predicate $F$ true of exactly the things that exist at worlds where everything happens as it actually does, the single axiom $\forall x\, Fx$ is maximally simple and maximally strong, and every truth becomes a law \citep{Lewis1983}.

\emph{Vague weighting.} Even with a fixed vocabulary, there is no precise exchange rate between simplicity and strength, and different rates may select different best systems.

\emph{Anthropocentrism.} Even if the rate were precise, it would reflect our standards of elegance, and laws of nature should not depend on what we find elegant \citep{vanFraassen1989}.

\section{Vocabulary}

The first version is answered by the restriction in Definition~2.2 of MET-12-a. Systems are stated in a vocabulary whose primitive predicates express genuine properties and relations, and MET-02 supplies the inventory: there are few genuine universals, and all other properties are constructed from them. The predicate $F$ expresses a property constructed, at enormous cost, from genuine ones, and the axiom $\forall x\, Fx$, rewritten in the genuine vocabulary, is as long as the description of the whole mosaic. Its apparent simplicity was an artefact of a primitive that does not express a genuine property. This is the proposal of \citet{Lewis1983}, and the theory of MET-02 gives it the inventory it needs.

\section{Vagueness in the weighting}

Let an \emph{admissible weighting} be an exchange rate between simplicity and strength consistent with the standards the sciences use in theory choice: one that prefers a system of a few short equations covering the whole mosaic to a list of its events, and prefers a list to a single tautology. The standards do not fix one weighting; they fix a range.

\begin{definition}\label{d:law}
A generalization is a \emph{law} if it is entailed by every best system on every admissible weighting.
\end{definition}

Definition~\ref{d:law} extends Definition~2.4 of MET-12-a, which already requires entailment by every best system, to quantify over admissible weightings as well. \citet{Lewis1994} held that if nature is kind, one system comes out best under any reasonable standards, and that if it is not, nothing deserves the name of law. Definition~\ref{d:law} makes the first half exact and refines the second.

\begin{proposition}\label{p:fewer}
Let $\mathcal{W}$ and $\mathcal{W}'$ be ranges of admissible weightings with $\mathcal{W} \subseteq \mathcal{W}'$. Every generalization that is a law relative to $\mathcal{W}'$ is a law relative to $\mathcal{W}$. In particular, the laws relative to a vague range are among the laws relative to each of its precise members.
\end{proposition}

\begin{proof}
A generalization is a law relative to $\mathcal{W}'$ if every best system on every weighting in $\mathcal{W}'$ entails it. Each weighting in $\mathcal{W}$ is in $\mathcal{W}'$, so every best system on every weighting in $\mathcal{W}$ entails it.
\end{proof}

Proposition~\ref{p:fewer} says that vagueness in the standards can only remove laws: widening the range of admissible weightings shrinks the set of laws, and never adds to it a generalization that some precise weighting would have excluded. Vagueness therefore yields fewer laws, not arbitrary ones. No generalization is a law merely because one weighting, rather than another equally admissible one, happens to be in force. The laws that survive are those that are robust across the range; in a world like ours, we expect the fundamental dynamical equations to be robust in this sense, since every admissible weighting prefers a system containing them both to a mere list of events and to a nearly empty system.

The definition is compatible with the epistemic account of vagueness of MET-18. On that account, \enquote{admissible} has a sharp but unknowable boundary, and the range of admissible weightings is a definite set whose edges we cannot locate. The laws are then a definite set, and Proposition~\ref{p:fewer} shows that our ignorance of the edges of the range induces ignorance only about generalizations that are robust across most admissible weightings and fail at the edges. For the generalizations that matter, which are robust across every weighting anyone would consider, there is no such ignorance.

\section{Anthropocentrism}

The third version says that the standards are ours. The reply is that the standards determine which summary of the mosaic counts as the best one, and that the mosaic is not ours. Two theorists with different standards, applied to the same mosaic, may disagree about which generalizations are laws only if the mosaic does not determine a system that is best by both; by Proposition~\ref{p:fewer}, the generalizations on which they disagree are not laws on the combined range. What depends on us is the concept of a law, as the concept of a summary depends on what a summary is for; what the concept picks out depends on the world. \citet{CohenCallender2009} draw the further conclusion that the choice of basic kinds may also be relative to a science; MET-12-a (Section~6) adopts that relativization for the special sciences and keeps the fundamental vocabulary fixed by the genuine properties.

\begin{objection}[Unkind worlds]
If the mosaic is such that every admissible weighting selects a different best system, Definition~\ref{d:law} counts no generalization as a law, and a world full of regularities has no laws.
\end{objection}

\begin{reply}
It counts as laws the generalizations common to all those systems, which in such a world will be few. We accept the consequence. A world whose regularities can be summarized in many equally good and incompatible ways is a world in which no single summary has the standing of law, and the intuition that it must nonetheless have laws is the governing conception, which MET-12-a rejects. Such a world still has regularities, and it still has causes and the dispositions that produce its events, which MET-12-a places in the mosaic and which do not depend on the laws.
\end{reply}

\section{Conclusion}

The problem of simplicity is three problems. Language relativity is removed by the vocabulary of genuine properties; the vagueness of the weighting is absorbed by requiring robustness across admissible weightings, so that vague standards yield fewer laws rather than arbitrary ones; and the dependence of the standards on us fixes only the concept of a best summary, not what the summary is of. The best system account survives each.

\begin{thebibliography}{9}\small
\bibitem[Cohen and Callender(2009)]{CohenCallender2009} Cohen, J. and Callender, C. (2009). A better best system account of lawhood. \emph{Philosophical Studies}, 145(1), 1--34.
\bibitem[Lewis(1983)]{Lewis1983} Lewis, D. (1983). New work for a theory of universals. \emph{Australasian Journal of Philosophy}, 61(4), 343--377.
\bibitem[Lewis(1994)]{Lewis1994} Lewis, D. (1994). Humean supervenience debugged. \emph{Mind}, 103(412), 473--490.
\bibitem[van Fraassen(1989)]{vanFraassen1989} van Fraassen, B.~C. (1989). \emph{Laws and Symmetry}. Oxford University Press.
\end{thebibliography}

\end{document}
